\begin{figure}[H]
\centering
\begin{tikzpicture}[>=Latex,x=1.02cm,y=0.82cm]
    \def\ang{20}
    \clip (-5.05,-0.65) rectangle (5.25,6.35);

    \fill[minkfieldred]
        (\ang:-5.8) -- (\ang:5.8) -- ++(90-\ang:7.1)
        -- ($(\ang:-5.8)+(90-\ang:7.1)$) -- cycle;

    \foreach \x in {-5,-4,...,5}
        \draw[mink grid x] (\x,-0.45) -- (\x,6.2);
    \foreach \y in {0,1,...,6}
        \draw[mink grid t] (-5.0,\y) -- (5.15,\y);
    \foreach \i in {-7,-6,...,7}{
        \draw[mink grid prime] (\ang:{\i*0.72}) -- ++(90-\ang:7.4);
        \draw[mink grid prime] (90-\ang:{\i*0.72}) -- ++(\ang:7.4);
    }

    \draw[mink dashed,gray!65] (-4.65,4.65) -- (4.65,-4.65);
    \draw[mink dashed,gray!65] (-4.65,-4.65) -- (4.65,4.65);
    \node[gray!65!black,font=\scriptsize,rotate=45] at (3.55,3.8)
        {línea de luz};

    \draw[mink axis] (-4.95,0) -- (5.1,0) node[below] {$x$};
    \draw[mink axis] (0,-0.5) -- (0,6.15) node[left] {$ct$};
    \draw[mink axis prime] (\ang:-5.0) -- (\ang:5.15) node[below right] {$x'$};
    \draw[mink axis prime] (90-\ang:-0.55) -- (90-\ang:6.3) node[right] {$ct'$};

    \coordinate (P) at (2.0,4.15);
    \fill[minkdarkred] (P) circle (2.8pt);
    \node[minkdarkred,above right=-1pt] at (P) {$P$};
    \draw[mink dashed,black!55] (P) -- (2.0,0)
        node[below,font=\scriptsize] {$x_P$};
    \draw[mink dashed,black!55] (P) -- (0,4.15)
        node[left,font=\scriptsize] {$ct_P$};
    \draw[mink dashed,minkdarkred] (P) -- ++(200:4.0);
    \draw[mink dashed,minkdarkred] (P) -- ++(250:4.0);

    \node[font=\small,align=left] at (-3.15,5.3)
        {\contour{white}{un evento, dos pares}\\
         \contour{white}{de coordenadas}};
\end{tikzpicture}
\caption{Diagrama de Minkowski con dos sistemas inerciales superpuestos. La rejilla azul corresponde a $S$ y la rejilla roja inclinada a $S'$. Las líneas de luz permanecen como fronteras comunes, mientras que el evento $P$ recibe coordenadas diferentes en cada marco.}
\label{fig:minkowski_original}
\end{figure}